\documentclass[10pt]{article}
\usepackage[utf8]{inputenc}
\usepackage[T2A]{fontenc}
\usepackage[ukrainian]{babel}
\usepackage{amsmath, amssymb, amsthm}
\usepackage{graphicx}
\usepackage{csquotes}
\usepackage{url}
\usepackage{geometry}
\geometry{a4paper, left=2cm, right=2cm, top=2cm, bottom=2cm}

\title{Теорема Нетер та стильниковий рушій Казимира: про не збереження імпульсу в асиметричних вакуумних системах}
\author{О. Дроздов\\
Харківський національний університет імені В. Н. Каразіна, майдан Свободи, 4, Харків, 61022, Україна\\
\texttt{alexey.drozdov.07@gmail.com}}
\date{}

\begin{document}

\maketitle

\begin{abstract}
У цій роботі аналізується стильниковий рушій Казимира через призму теореми Нетер. Розглянута система — ідеально провідна пластина з нанорозмірною сотовою структурою з одного боку — не має трансляційної симетрії вздовж осі, нормальної до пластини, через геометричну асиметрію вакуумних мод з двох її сторін. Як наслідок, гамільтоніан вакууму явно залежить від координати $l$ (ефективної довжини резонатора), і, згідно з теоремою Нетер, відповідна компонента імпульсу не зберігається. Це не збереження проявляється як результуюча сила квантового вакууму (тяга Казимира). Показано, що отриманий вираз для сили на одиницю площі безпосередньо відображає цю порушену симетрію й узгоджується як з гамільтоновим підходом, так і з різницею густин енергії. Результат підкреслює глибокий фізичний принцип: макроскопічна тяга від квантових флуктуацій вакууму виникає саме тому, що система не є інваріантною відносно просторових зсувів.
\end{abstract}

\medskip
\noindent\textbf{Ключові слова:} теорема Нетер, ефект Казимира, квантовий вакуум, не збереження імпульсу, наносоти, тяга Казимира

\medskip
\noindent\textbf{УДК:} 53.09, 621.3, 629.7 \quad
\textbf{Номери PACS:} 11.10.Ef, 03.70.+k, 42.50.Lc

\section{Вступ}
Теорема Нетер встановлює глибокий зв’язок між неперервними симетріями фізичної системи та збереженими величинами. Зокрема, інваріантність відносно просторових зсувів означає збереження імпульсу. У класичній конфігурації Казимира з двома пластинами система симетрична в площині пластин, тому поперечний імпульс зберігається; лише нормальна компонента впливається умовами на межі, але й тут сили врівноважені між двома тілами.

Навпаки, \textit{однотіловий} стильниковий рушій Казимира, запропонований у роботі \cite{Drozdov2024}, явно порушує трансляційну симетрію вздовж осі $z$ через різні граничні умови з двох боків єдиної пластини: з одного боку — гладка поверхня, з іншого — наноструктуровані комірки. Ця асиметрія робить вакуумний гамільтоніан $\langle 0|\hat{\mathcal{H}}|0\rangle$ функцією геометричного параметра $l$.

\section{Теорема Нетер і не збереження імпульсу}
Розглянемо вакуумний гамільтоніан на одиницю площі, отриманий у \cite{Drozdov2024}:
\begin{equation}
\frac{\langle 0|\hat{\mathcal{H}}|0\rangle}{S} = \frac{\hbar c}{a^2\pi} \Bigg\{
l \sum_{n_x=(0)1}^{\infty}\sum_{n_y=(0)1}^{\infty}
\int_0^{\infty} \sqrt{ \frac{\pi^2}{a^2}(n_x^2 + n_y^2) + k_z^2 } \, dk_z
+ (L - l) \int_0^{\infty}\!\!\int_0^{\infty}
\int_0^{\infty} \sqrt{k_x^2 + k_y^2 + k_z^2} \, dk_z \, \frac{a}{\pi}dk_x \frac{a}{\pi}dk_y
\Bigg\}.
\end{equation}
Цей вираз явно залежить від $l$. Отже, система \textbf{не інваріантна} відносно зсувів $z \to z + \delta z$, бо таке перетворення змінює ефективну геометрію резонаторів зі структурованого боку.

Згідно з теоремою Нетер, генератор просторових зсувів — компонента імпульсу $P_z$ — зберігається \textbf{лише тоді}, коли дія (або гамільтоніан у статичному випадку) не залежить від $z$. Оскільки $\partial \langle H \rangle / \partial l \neq 0$, симетрія порушена, і $P_z$ не зберігається.

Відповідна сила:
\begin{equation}
\frac{F}{S} = \frac{\partial}{\partial l} \frac{\langle 0|\hat{\mathcal{H}}|0\rangle}{S}
= \frac{\hbar c}{a^2\pi} \left\{
\sum_{n_x=(0)1}^{\infty}\sum_{n_y=(0)1}^{\infty}
\int_0^{\infty} \sqrt{ \frac{\pi^2}{a^2}(n_x^2 + n_y^2) + k_z^2 } \, dk_z
- \int_0^{\infty}\!\!\int_0^{\infty}
\int_0^{\infty} \sqrt{k_x^2 + k_y^2 + k_z^2} \, dk_z \, dn_x dn_y
\right\},
\end{equation}
що збігається з тягою Казимира, отриманою в \cite{Drozdov2024}.

Отже, поява чистої сили — не парадокс, а передбачуваний наслідок порушення трансляційної симетрії. Квантовий вакуум «штовхає» сильніше з одного боку через різну структуру мод, і теорема Нетер пояснює, чому імпульс не має зберігатися в такій системі.

\section{Висновки}
Стильниковий рушій Казимира є яскравим прикладом системи, де теорема Нетер прямо пояснює походження макроскопічної сили: \textbf{відсутність просторової однорідності призводить до не збереження імпульсу}, а отже — до тяги. Це розуміння підкреслює, що ефекти квантового вакууму є не просто математичними артефактами, а проявами глибоких принципів симетрії. Експериментальне підтвердження цієї тяги стане не лише доказом існування віртуальних фотонів, а й перевіркою зв’язку між геометрією, симетрією та динамікою в квантовій теорії поля.

\begin{thebibliography}{9}
\bibitem{Drozdov2024}
О. Дроздов, \textit{Стильниковий рушій Казимира: про силу, що діє на соти на пластині, які ідеально проводять}, Вісник Харківського національного університету імені В. Н. Каразіна. Серія «Фізика», вип. 41, 28--41 (2024). \url{https://doi.org/10.26565/2222-5617-2024-41-04}
\bibitem{Noether1918}
Е. Нетер, \textit{Invariante Variationsprobleme}, Nachr. d. König. Gesellsch. d. Wiss. zu Göttingen, Math-phys. Klasse, 235–257 (1918).
\bibitem{Casimir1948}
Х. Б. Г. Казимир, \textit{Про притягання між двома ідеально провідними пластинами}, Proc. K. Ned. Akad. Wet. B \textbf{51}, 793 (1948).
\end{thebibliography}

\end{document}
